\documentclass[10pt,a4paper]{article}
\usepackage[T1]{fontenc}
\usepackage[utf8]{inputenc}
\usepackage[margin=16mm]{geometry}
\usepackage{booktabs}
\usepackage[table]{xcolor}
\usepackage{microtype}
\usepackage{lmodern}
\usepackage[hidelinks]{hyperref}

\definecolor{ink}{HTML}{193A52}
\definecolor{stripe}{HTML}{EDF3F6}
\setlength{\parindent}{0pt}
\setlength{\parskip}{3pt}
\setlength{\tabcolsep}{7pt}
\renewcommand{\arraystretch}{1.12}
\pagestyle{empty}
\newcommand{\reportsection}[1]{\vspace{2mm}{\large\bfseries\color{ink}#1}\par\vspace{0.5mm}}

\begin{document}
{\LARGE\bfseries\color{ink} Lab 03: Chaotic Text Generator}\par
\vspace{1mm}
{\normalsize\textbf{Author: Ayoub El Houani} \hfill One-page project report}\par
\vspace{1mm}\hrule\vspace{1mm}

\reportsection{1. Objective and method}
This project compares a plain part-of-speech HMM, a role-split HMM and a trigram generator. A Streamlit application generates batches, displays tags, computes metrics, exports CSV and tags user input with log-space Viterbi. The synthetic corpus is split by document into 960 training and 240 test sentences. Its small grammar makes this a teaching experiment, not a measure of real English quality.

\reportsection{2. Generation results at $T=1.0$}
Each model generated 500 sentences with seed 42. \textit{Valid} means the tag pattern was seen in training; \textit{plausible} means the first noun--verb pair was seen; \textit{novel} means the full sentence was not copied; distinct-2 is the fraction of unique word bigrams.
\begin{center}
\rowcolors{2}{stripe}{white}
\begin{tabular}{lrrrr}
\toprule
\textbf{Model} & \textbf{Valid pattern} & \textbf{Plausible pair} & \textbf{Novel} & \textbf{Distinct-2} \\
\midrule
Plain HMM      & 16.2\%  & 14.8\% & 99.8\% & 0.198 \\
Role-split HMM & 100.0\% & 60.6\% & 99.4\% & 0.148 \\
Trigram        & 57.0\%  & 76.0\% & 91.2\% & 0.105 \\
\bottomrule
\end{tabular}
\end{center}

\textbf{Interpretation.} The plain HMM reuses POS states at different sentence positions, so it can combine incompatible transitions. Role-split states retain subject, object or prepositional roles. They produce 100\% valid patterns in this sample, but do not learn meaning or guarantee grammatical English. Trigram tags are estimated with the role-split HMM.

\reportsection{3. Settings and generated examples}
For a \textbf{writing toy}, $T=1.5$ is a reasonable choice for variety; for \textbf{autocomplete}, $T=0.5$ favours common continuations. These are heuristics, not measured autocomplete results.

With the role-split HMM, seed 42 and $T=1$, ``the children eat the bread at the park .'' is a plausible novel sentence. ``the farmers eat the river by the school .'' has a valid structure but fails in \textbf{meaning}: independent word emissions do not model semantic compatibility. Its noun--verb pair was seen, showing the limits of the plausibility proxy.

\reportsection{4. Evaluation on real text}
Both HMMs achieved 100\% token accuracy on the 240 held-out synthetic sentences. On \href{https://github.com/UniversalDependencies/UD_English-EWT}{UD English EWT} (12,544 training and 2,077 test sentences; CC BY-SA 4.0), the measured results were:
\begin{center}
\rowcolors{2}{stripe}{white}
\begin{tabular}{lrrrr}
\toprule
\textbf{Tagger} & \textbf{Accuracy} & \textbf{Known} & \textbf{Unknown} & \textbf{Whole sentence} \\
\midrule
Most-frequent tag & 86.2\% & 90.8\% & 28.6\% & 31.9\% \\
HMM + Viterbi     & 88.0\% & 92.9\% & 28.2\% & 40.4\% \\
\bottomrule
\end{tabular}
\end{center}

Unknown words make up 7.5\% of test tokens. Real web text has 17 tags, more varied patterns, ambiguity and unseen words; the model's smoothed emissions lack word-shape information for unknowns. The synthetic grammar is much smaller and more regular. Results can vary with treebank revisions.
\end{document}
